\section{The Trained Continuation as the Certified Bellman Object}
\label{sec:r41direct}
A learned continuation enters the economic argument only when the function
actually returned by training, rather than its training target, satisfies the
Bellman comparison. We now set $v_t=f_t$ in
Theorem~\ref{thm:r39comparison}, where $f_0,\ldots,f_T$ are the stored neural
networks. The continuation used inside the date-$t$ Bellman operator is
$f_{t+1}$. Its replacement by a reference approximation is not permitted in
this calculation. This construction completes the connection between the
learned scalar continuation and the returned dynamic policy; it does not
require the training procedure to have found a global minimizer.

\subsection{A certificate that does not use the training reference}
Let $\{x_i\}$ and $\{a_j\}$ cover $X$ and $A$ with radii $h_x$ and $h_a$,
respectively. Ties in nearest-node selection have a fixed measurable rule.
Suppose $f_t$ has Lipschitz constant $L_t$, and let $L_{x,t}$ and $L_{a,t}$
bound the variation of $Q_t(f_{t+1};x,a)$ in state and action. Verified
arithmetic returns
\[
 f^-_{ti}\leq f_t(x_i)\leq f^+_{ti},\qquad
 q^-_{tij}\leq Q_t(f_{t+1};x_i,a_j)\leq q^+_{tij}.
\]
These intervals enclose evaluations of the network coefficients and the
primitive transition and cost. Define $L_{ti}=\min_jq^-_{tij}$ and
$U_{ti}=\min_jq^+_{tij}$. Store any feasible grid action $a_{j(t,i)}$;
minimizing $q^+_{tij}$ is one implementable choice.

\begin{proposition}[Direct neural Bellman certificate]
\label{prop:r41direct}
Set
\begin{align}
 \ell_t&=\min_i(f^-_{ti}-U_{ti})-(L_t+L_{x,t})h_x,
       \label{eq:r41direct}\\
 u_t&=\max_i(f^+_{ti}-L_{ti}+L_{a,t}h_a)
                         +(L_t+L_{x,t})h_x,\nonumber\\
 \eta_t&=\max_i(q^+_{ti,j(t,i)}-L_{ti})+L_{a,t}h_a+2L_{x,t}h_x.
       \nonumber
\end{align}
If $\ell_T\leq f_T-g\leq u_T$ on $X$, the nearest-node actor $\pi$ satisfies
\begin{equation}
 0\leq J_0^\pi(x)-V_0^*(x)\leq
 E_f:=\sum_{t<T}B_t(u_t-\ell_t+\eta_t)+B_T(u_T-\ell_T).
 \label{eq:r41E}
\end{equation}
Its comparison class is all feasible adapted policies. If a separate
own-policy calculation gives $J_0^\pi(x)\leq\overline J_0^\pi(x)$, then
\begin{equation}
 J_0^\pi(x)-V_0^*(x)
 \leq\overline J_0^\pi(x)-f^-_0(x)+\sum_{t=0}^TB_tu_t.
 \label{eq:r41query}
\end{equation}
Neither certificate reads a reference spline, its nodal values, or its
Bellman brackets.
\end{proposition}
\begin{proof}
At a node the continuous-action minimum lies in
$[L_{ti}-L_{a,t}h_a,U_{ti}]$: every feasible minimizer is within $h_a$ of a
grid action. Subtraction gives the nodal residual interval. The difference
$f_t-T_tf_{t+1}$ has Lipschitz constant at most $L_t+L_{x,t}$, which proves
the first two bounds throughout each covering cell. Moving a chosen action
and the Bellman minimum from a node to the same state costs at most
$2L_{x,t}h_x$. This proves the actor bound. Apply
Theorem~\ref{thm:r39comparison}, whose lower value bound also proves
\eqref{eq:r41query}. The argument uses the stored neural functions at both
successive dates, including the terminal date.
\end{proof}

In contrast to \eqref{eq:r39mesh}, the state remainder includes $L_t$:
the trained network need not interpolate the verification nodes. Omitting
that term would incorrectly treat the training target as the learned
function. Additive shifts $f_t\mapsto f_t+k_t$ change both residual endpoints
by $k_t-\beta_tk_{t+1}$ and leave their oscillation and the actor allowance
unchanged. The terminal shift completes the telescoping sum. Thus both
\eqref{eq:r41E} and the exact version of \eqref{eq:r41query} are invariant
to the arbitrary levels of the learned continuations.

For the scalar economy, a stored ReLU network has rational breakpoints and
rational affine coefficients because its floating coefficients are exact
dyadic numbers when interpreted as stored data. On each activation region
we compute its affine coefficients from those weights. This is an exact
compilation of the trained network, not a fitted spline or a conventional
Bellman solution. If a breakpoint is rounded by at most $e_r$, choosing an
adjacent affine extension adds at most $2L_te_r$ to the value enclosure.
Rational coefficient enclosures and the interval width of the evaluation
argument are charged separately. Appendix~\ref{app:r41arithmetic} proves
this statement and specifies the arithmetic checks.

The terminal cost in \eqref{eq:r39capital} is piecewise quadratic. Hence the
extrema of $f_T-g$ occur at an activation breakpoint, at $3/8$, at a domain
endpoint, or at a stationary point in an intervening region. All these
points are calculated in rational arithmetic. Terminal fit error is thereby
included in the same neural chain, rather than set to zero because the
terminal primitive is known.

The deposited networks were trained on the earlier conventional nodal
values. That training-data cost remains part of their provenance. The present
verification changes the certified mathematical object, not that history:
there are no new neural training runs, and there is no assertion that the
reference was unnecessary for the original fit. The verification program
continues to run when the conventional spline evaluator is disabled.

\subsection{From the stored actor to its numerical execution}
\label{sec:r41deployment}
The stored actor is not the last object entering an economic calculation.
Let the observed state be rounded with error at most $s_x$. Suppose selected
actions and grid cutpoints are represented exactly, and their comparisons
are exact. Nearest-node selection at the rounded state can choose a different
cell, but its node is within $h_x+s_x$ of the true state. No continuity of
the tabulated actor is required. Suppose further that the numerical
transition and stage-cost graphs satisfy
\[
 \|\widetilde F_t(x,a,z)-F_t(x,a,z)\|\leq s_{F,t},\quad
 |\widetilde c_t(x,a)-c_t(x,a)|\leq s_{c,t},\quad
 |\widetilde g(x)-g(x)|\leq s_g.
\]
The bounds include conversion of their inputs and clipping. Both transitions
stay in $X$. The recursive risk functional itself retains its stated
mathematical meaning and is evaluated by enclosing intervals.

\begin{proposition}[Two execution contracts]
\label{prop:r41deployment}
Under Proposition~\ref{prop:r41direct}, the rounded-state actor
$\pi^{\rm exec}$ applied to the original economic transition satisfies
\begin{equation}
 0\leq J_0^{\pi^{\rm exec}}-V_0^*
       \leq E_f+\sum_{t<T}2B_tL_{x,t}s_x.
 \label{eq:r41actor_execution}
\end{equation}
For the numerical economy that also uses $\widetilde F_t$,
$\widetilde c_t$, and $\widetilde g$, its recursive cost instead satisfies
\begin{equation}
 \widetilde J_0^{\pi^{\rm exec}}-V_0^*
 \leq E_f+E_{\rm exec},\quad
 E_{\rm exec}:=\sum_{t<T}B_t
 (2L_{x,t}s_x+s_{c,t}+\beta_tL_{t+1}s_{F,t})+B_Ts_g.
 \label{eq:r41execution}
\end{equation}
The left side of \eqref{eq:r41execution} need not be nonnegative: its two
values belong to different transition and accounting models.
\end{proposition}
\begin{proof}
The nearest-node proof of Proposition~\ref{prop:r41direct} replaces $h_x$
by $h_x+s_x$ in the actor allowance. Feasibility is unchanged. This proves
\eqref{eq:r41actor_execution}. Monotonicity and cash invariance give
\[
 |\rho_t(f_{t+1}(\widetilde F_t))-
   \rho_t(f_{t+1}(F_t))|\leq L_{t+1}s_{F,t}.
\]
Thus the numerical chosen-action operator has an additional upper error
$2L_{x,t}s_x+s_{c,t}+\beta_tL_{t+1}s_{F,t}$ relative to the original Bellman
operator. Backward induction on the numerical policy, with terminal error
$s_g$, adds the discounted sum of these errors to the proof of
Theorem~\ref{thm:r39comparison}. The lower bound on the original optimum is
unchanged. This proves \eqref{eq:r41execution} without identifying a rounded
state trajectory with an exact one.
\end{proof}

We instantiate the second contract with separately rounded binary32
operations, round-to-nearest, and gradual underflow. On the specified domain
there is no overflow. With $u=2^{-24}$ and $\alpha=2^{-150}$, each elementary
rounding obeys $|\operatorname{fl}(y)-y|\leq u|y|+\alpha$. Rational
range-and-error propagation through the actual arithmetic graph gives
\[
 s_x<5.961\,10^{-8},\quad s_F<3.167\,10^{-7},\quad
 s_c<5.267\,10^{-7},\quad s_g<6.413\,10^{-7}.
\]
The stage-cost bound displayed here covers both prices. All action values
and cutpoints at the reported dyadic resolutions are exactly representable;
this accounts for zero action-storage error without setting the whole
execution allowance to zero. Clipping is nonexpansive. Fused evaluation,
flush-to-zero, altered state acquisition, and a physical actuator would
require their own contracts. The software experiment evaluates the rounded
transition and rounded costs, not a claim of physical deployment validation.
